\def\topic{}
\ifdefined\MyHeaderIncluded
\else
\def\qwertyflag{true}

\ifdefined\MyHeaderIncluded

\else
  \newcommand{\foobar}{FooBar}
  \def\MyHeaderIncluded{true}
% if is till the end of the file

\documentclass[10pt,a4paper]{article}
\usepackage[utf8]{inputenc}
%\usepackage[german]{babel}
%\usepackage[cyrillic]{babel}
\usepackage[T1]{fontenc}
%\usepackage[T2A]{fontenc}
%\usepackage[english, russian]{babel}
\usepackage{amsmath}
\usepackage{amsfonts}
\usepackage{amssymb}
\usepackage{amsthm}
\usepackage{mathtools} %coloneqq

\usepackage{bussproofs} %

%\renewcommand{\thesubsection}{\arabic{subsection}}
%\newtheorem{thm}{Theorem}[section]
%\renewcommand{\thethm}{\arabic{subsection}.\arabic{thm}}

\usepackage{thmtools}
\usepackage{hyperref}
\hypersetup{
    colorlinks=true,
    linkcolor=blue,
    filecolor=magenta,      
    urlcolor=cyan,
    pdftitle={Overleaf Example},
    %pdfpagemode=FullScreen,
    }
%\usepackage{stix-type1}%\turnediota
\theoremstyle{definition}

\newtheorem{constructor}{Constructor}[section]
\newtheorem{iaxiom}{Axiom}
\newtheorem{axiom}{Axiom}
\newtheorem*{axiom*}{Axiom}
\newtheorem{definition}{Definition}[section]
\newtheorem*{definition*}{Definition}
\newtheorem{remark}{Remark}[section]
\newtheorem*{remark*}{Remark}
\newtheorem{recall}{Recall}[section]
\newtheorem*{recall*}{Recall}

\newtheorem{lemma}{Lemma}[section]
\newtheorem{theorem}{Theorem}[section]
\newtheorem*{theorem*}{Theorem}
\newtheorem{note}{Note}[section]
\usepackage{etoolbox}

%\newcommand*{\@Dashv}[2]{%
%  \reflectbox{$\m@th#1#2$}%
%}
\newcommand*{\separation}[3]{% 1)binder, 2)set, 3)condition
  \{#1\in #2: #3\}
}

% https://www.freecodecamp.org/news/simple-regex-tricks-for-beginners-3acb3fa257cb/
%  \\mathcal\{T\} |---> \topology
%  \\mathcal\{([^\}]*)\}   |--->  $1
\newcommand{\AND}{\ \ \&\ \ }
\newcommand{\any}{\mathrm{\ any}}
\newcommand{\borel}{\mathcal{B}}
\newcommand{\can}{\mathrm{can}}
\newcommand{\canonize}{\mathrm{canonize}}
\newcommand{\closure}{\mbox{cl}}
\newcommand{\complex}{\mathbb{C}}
\newcommand{\converges}{\longrightarrow}
\newcommand{\countable}{\mbox{\ countable}}
\newcommand{\ldefined}{\mbox{\ defined}}
\newcommand{\dom}{\mbox{dom}}
\newcommand{\Dom}{\mbox{Dom}}
\newcommand{\Done}{\noindent \textbf{Done.}\\}
\newcommand{\lexists}{\mbox{\ exists}}
\newcommand{\lforall}{\mbox{\ forall}}
\newcommand{\lunique}{\mbox{\ unique}}
\newcommand{\family}{\mathrm{\ family}}
\newcommand{\Fm}{\mathrm{Fm}}
\newcommand{\FuncSymb}{\mathrm{FuncSymb}}
\newcommand{\function}{\mathrm{\ function}}
\newcommand{\functional}{\mbox{\ functional}}
\newcommand{\IFF}{\ \Longleftrightarrow\ }
\renewcommand{\iff}{\ \longleftrightarrow\ }
\newcommand{\IMPLIES}{\ \Longrightarrow\ }
\renewcommand{\implies}{\ \longrightarrow\ }
\newcommand{\Inf}{\mathrm{Inf}}
\newcommand{\injective}{\mbox{\ injective}}
\newcommand{\injection}{\mbox{\ injection}}
\newcommand{\interior}{\mbox{int}}
\newcommand{\intersection}{\bigcap}
\newcommand{\isempty}{\mathrm{\ empty}}
\newcommand{\LB}{\mathrm{LB}}
\newcommand{\nghbset}{\mathcal{N}}
\newcommand{\nonempty}{\mathrm{\ nonempty}}
\newcommand{\Ord}{\mathrm{Ord}}
\newcommand{\power}{\mathcal{P}}
\newcommand{\pr}{\mathrm{pr}}
\newcommand{\fst}{\mathrm{fst}}
\newcommand{\snd}{\mathrm{snd}}
\newcommand{\PredSymb}{\mathrm{PredSymb}}
\newcommand{\Proof}{\noindent \textbf{Proof.}\\}
\newcommand{\ran}{\mathrm{ran}}
\newcommand{\Ran}{\mathrm{Ran}}
\newcommand{\real}{\mathbb{R}}
\newcommand{\relation}{\mathrm{\ relation}}
\newcommand{\set}{\mathrm{\ set}}
\newcommand{\singleton}{\mathrm{\ singleton}}
\newcommand{\semantics}{\mathrm{semantics}}
\newcommand{\semialgebra}{\mathcal{S}}
\newcommand{\Sup}{\mathrm{Sup}}
\newcommand{\surjective}{\mbox{\ surjective\ on\ }}
\newcommand{\Term}{\mathrm{Term}}
\newcommand{\Type}{\mathrm{Type}}
\newcommand{\topology}{\mathcal{T}}
\newcommand{\total}{\mbox{\ total\ on\ }}
\newcommand{\transitive}{\mathrm{\ transitive}}
\newcommand{\union}{\bigcup}
\newcommand{\lTrue}{\mathrm{True}}
\newcommand{\lFalse}{\mathrm{False}}
\newcommand{\lUncertain}{\mathrm{Uncertain}}
\newcommand{\lif}{\mathrm{if}}
\newcommand{\lthen}{\mathrm{then}}
\newcommand{\lelse}{\mathrm{else}}
\newcommand{\lexception}{\mathrm{exception}}
\newcommand{\lundef}{\mathrm{undef}}
\newcommand{\univocal}{\mathrm{\ univocal}}
\newcommand{\univocalat}{\mathrm{\ univocal\ at\ }}
\newcommand{\UB}{\mathrm{UB}}
\newcommand{\Un}{\mathcal{U}}
\newcommand{\Var}{\mathrm{Var}}

\newcommand{\notvarphi}{\phi}
\renewcommand{\phi}{\varphi}

\newcommand{\comment}[1]{\quad\mbox{(#1)}}
% local commands
%\preto\section{\undef\blipo} %etoolbox
\preto\subsection{\undef\blipo} %etoolbox
% https://tex.stackexchange.com/questions/105684/limit-macro-newcommand-scope-to-section
\usepackage{graphicx}
\usepackage{authblk}

\setlength\parindent{0pt} % whole doc of \noindent
\usepackage{soul}
%\usepackage{mathgl}

\fi % guard
\ifdefined\topic
\else
\def\topic{No Titel}
\fi
\title{\topic}
\author[1]{Georgy Dunaev}
\begin{document}
\maketitle
\fi
\section{Restricted quantifiers: notation and a syntactic proposal}
For an explicitly specified bound variable $x$ that is not free in $A$,
restricted quantifiers abbreviate ordinary quantification with a condition:
\[
(\forall x\in A,\ Q)\iff(\forall x,\ x\in A\implies Q),
\qquad
(\exists x\in A,\ Q)\iff(\exists x,\ x\in A\land Q).
\]
Likewise, $(!x\in A,\ Q)$ means $(!x,\ x\in A\land Q)$, with $!$
expressing at most one witness. All substitutions are capture-avoiding.

The commented draft below explores a more compressed syntax in which a
bound variable would be recovered from the first term of the restriction.
That proposal needs an explicit rule for selecting the variable and
handling terms with several free variables. Until such a rule is fixed,
the explicit-binder notation above determines the meaning unambiguously.
\subsection{A precise version of inferred binders}
One way to make the proposed compressed notation unambiguous is to
supply a syntactic parameter context $P$: a finite set of variable names
that are parameters, not candidates for the new binder. For a restriction
formula $R$, compute $\mathrm{FV}(R)\setminus P$. Permit the abbreviation
only when this set is a singleton $\{x\}$. Then expand
\[
\forall[R]\,Q\quad\text{as}\quad\forall x\,(R\implies Q),
\qquad
\exists[R]\,Q\quad\text{as}\quad\exists x\,(R\land Q).
\]
The analogous at-most-one abbreviation expands to $!x\,(R\land Q)$.
When there are no candidates or several, require an explicit binder;
do not choose a variable from its alphabetical order or an accidental
presentation order. This is one formal realization of the proposal,
not a claim that the current checker implements it.

\begin{proposition}[Hygiene of the elaborated restriction]
Once the binder has been determined, capture-avoiding substitution for
parameters commutes with the displayed expansion, provided the binder
is first renamed fresh for all replacement terms. Bijective renaming
of variable names, applied also to $P$, preserves the inference rule.
\end{proposition}
\Proof
For the universal abbreviation, expansion yields $\forall x(R\implies Q)$.
Freshness makes its substitution exactly
$\forall x(R\sigma\implies Q\sigma)$, which is the expansion with the
same explicitly retained binder and substituted components. Existential
and at-most-one cases use the identical argument with conjunction.
A bijective renaming $\rho$ sends the candidate set to
$\rho(\mathrm{FV}(R)\setminus P)$, because free occurrences and parameter
names are renamed together. A singleton remains the singleton
$\{\rho(x)\}$, so the inferred binder is renamed consistently.
\Done

\begin{example}[Inference precedes substitution]
With $P=\{A\}$, the restriction $x\in A$ has the unique candidate $x$.
If $A$ is replaced by the variable $y$ while the parameter context is
left unchanged, the unelaborated restriction $x\in y$ has two candidate
names. Repeating inference would reject it or select a different binder
under an arbitrary heuristic. Elaborating first retains the intended
binder $x$, and ordinary capture-avoiding substitution gives
$\forall x(x\in y\implies Q\sigma)$. A parameter-context update could
also record $y$ explicitly; it must not be implicit in the mathematical
meaning of an already parsed expression.
\end{example}

%\section{Restricted quantifiers}
%TODO:\\
%The first formula is the ... .\\
%The first term in formula is the leftmost term in the first primitive formula.\\
%Here $R_x$ means that variable x is the first free variable in the proposition $R$.\\

%\begin{constructor} Restricted quantifier.\\
%If $R_x$ is a proposition with free variable x,\\
%\quad $Q$ is a proposition,\\
%\quad $\#$ is one of quantifiers ($\forall$, $\exists$ or $!$)\\
%then $\# R:Q$ is a formula.
%\end{constructor}

%\begin{definition} Restricted forall.\\
%\label{def:restricted forall}
%Meta--context: assume that the first term of a proposition R is a free %variable x.
%\[
%(\forall R:Q)
%\iff
%(\forall x: R \implies Q)
%\]
%\end{definition}

%\begin{definition} Restricted forall.\\
%\label{def:restricted forall}
%Meta--context 0: assume that the first term of a proposition $R$ is $t$, %and it has no free variables.
%\[
%(\forall R:Q)
%\iff
%(R \implies Q)
%\]
%Meta--context 1: assume that the first term of a proposition $R$ is $t$, and the first free variable of $t$ is $x$.
%\[
%(\forall R:Q)
%\iff
%(\forall x: R \implies Q)
%\]
%\end{definition}


%\begin{definition} Restricted exists.\\
%\label{def:exists restricted }
%Meta--context: assume that the first term of a proposition R is a free variable x.
%\[
%(\exists R:Q)
%\iff
%(\exists x: R \land Q)
%\]
%\end{definition}

%\begin{definition} Restricted unique.\\
%\label{def:unique restricted }
%Meta--context: assume that the first term of a proposition R is a free %variable x.
%\[
%(! R:Q)
%\iff
%(! x: R \land Q)
%\]
%\end{definition}

\ifdefined\qwertyflag
\undef\qwertyflag
\end{document}
\fi
